\documentclass[11pt]{ctexart}
\usepackage[a4paper,top=2cm,bottom=2cm,left=2.2cm,right=2.2cm]{geometry}
\usepackage{fancyvrb}
\usepackage{xcolor}
\usepackage{colortbl}
\usepackage{fontspec}
\usepackage[most]{tcolorbox}
\usepackage{booktabs}
\usepackage{enumitem}
\usepackage{titlesec}
\usepackage{hyperref}

\setmonofont{Consolas}
\setCJKmonofont{SimHei}
\setlength{\parindent}{0pt}
\setlist{nosep,leftmargin=1.4em}

% ── Step4 主题色（与 Step1/2/3 同族、换为深绛红） ──
\definecolor{deep}{RGB}{96,32,44}        % 主色：深绛红
\definecolor{accent}{RGB}{170,52,64}     % 强调：绛红
\definecolor{keyline}{RGB}{90,130,200}   % 知识点：蓝
\definecolor{keybg}{RGB}{238,244,255}
\definecolor{warnline}{RGB}{200,120,60}  % 警告：橙
\definecolor{warnbg}{RGB}{255,243,235}
\definecolor{okline}{RGB}{26,116,88}     % 成果：绿
\definecolor{okbg}{RGB}{240,252,244}
\definecolor{tipline}{RGB}{180,110,120}  % 提示：绛
\definecolor{tipbg}{RGB}{250,240,242}
\definecolor{codebg}{RGB}{249,249,252}   % 代码底
\definecolor{rowtint}{RGB}{248,238,240}  % 表头浅绛

\titleformat{\section}{\large\bfseries\color{deep}}{\thesection}{0.6em}{}
\titleformat{\subsection}{\normalsize\bfseries\color{accent}}{\thesubsection}{0.5em}{}
\titleformat{\part}[display]{\centering\Huge\bfseries\color{deep}}{}{0pt}{\Huge}
\titlespacing{\section}{0pt}{14pt}{6pt}
\titlespacing{\subsection}{0pt}{10pt}{4pt}

\newtcolorbox{keybox}[1]{enhanced, colback=keybg, colframe=keyline,
  boxrule=0.6pt, arc=1.5mm, left=3mm, right=3mm, top=2mm, bottom=2mm,
  title={\bfseries #1}, fonttitle=\small\bfseries, coltitle=white,
  colbacktitle=keyline}
\newtcolorbox{warnbox}[1]{enhanced, colback=warnbg, colframe=warnline,
  boxrule=0.6pt, arc=1.5mm, left=3mm, right=3mm, top=2mm, bottom=2mm,
  title={\bfseries #1}, fonttitle=\small\bfseries, coltitle=white,
  colbacktitle=warnline}
\newtcolorbox{okbox}[1]{enhanced, colback=okbg, colframe=okline,
  boxrule=0.6pt, arc=1.5mm, left=3mm, right=3mm, top=2mm, bottom=2mm,
  title={\bfseries #1}, fonttitle=\small\bfseries, coltitle=white,
  colbacktitle=okline}
\newtcolorbox{tipbox}[1]{enhanced, colback=tipbg, colframe=tipline,
  boxrule=0.5pt, arc=1.5mm, left=3mm, right=3mm, top=1.5mm, bottom=1.5mm,
  title={\bfseries #1}, fonttitle=\small\bfseries, coltitle=white,
  colbacktitle=tipline}

\newcommand{\code}[1]{\texttt{#1}}
\DefineVerbatimEnvironment{Verilog}{Verbatim}{frame=single,fontsize=\small}

\begin{document}

% ==================== 封面 ====================
\begin{titlepage}
\centering
\vspace*{2.2cm}
{\color{deep}\Huge\bfseries OSD 与系统联调}\\[6pt]
{\color{deep}\Large\bfseries 从零手搓 Step4 教程}\\[14pt]
{\large 2026 FPGA 创新设计赛道（安路）赛题二}\\[6pt]
{\large 视觉方向 · Step4\_osd\_integration}\\[20pt]

\begin{tcolorbox}[enhanced, colback=codebg, colframe=accent, boxrule=0.8pt,
  arc=2mm, width=12.5cm, left=4mm, right=4mm, top=3mm, bottom=3mm]
\centering
{\small 读者：零基础队员（建议先读完 Step1～3 教程）\\
目标：把检测结果画回画面——文字/准星/目标框，并完成整机联调\\
前置要求：会 Verilog、懂 Step2 检测链与 Step3 闭环的接口\\
预计学习时间：10～14 天（每天 2～3 小时）}
\end{tcolorbox}

\vspace{8mm}
\begin{tcolorbox}[enhanced, colback=tipbg, colframe=tipline, boxrule=0.5pt,
  arc=2mm, width=12.5cm, left=4mm, right=4mm, top=2mm, bottom=2mm]
\centering
{\small\color{deep} \textbf{本教程与你手上的代码是什么关系}\\
Step2 在"看"、Step3 在"动"；Step4 把这一切\textbf{画给你看}——\\
左上角状态行、绿色准星、红色目标框。评委看到的画面从这里来。}
\end{tcolorbox}

\vfill
{\small\color{black!50} fpga-ptz-tracker 教程系列 · 2026-09}
\end{titlepage}

\tableofcontents
\newpage

% ============================================================
\part{从零学习教程}

% ============================================================
\section{你在造什么}

\subsection{一句话}

\textbf{在视频画面上叠一层"仪表盘"：系统状态、目标十字、包围框，}
\textbf{再把检测、云台、显示拼成一台完整的机器。}

\subsection{整机长什么样}

\begin{Verilog}
                    +--------------------------------------------+
 96bit 视频流 --->  |  直通出（原样给显示链路，零扰动）            |
        |           |                                            |
        +--> detect_top（检测）
                  |----> 135bit 目标包 --CDC--> osd_track --+--> OSD 叠加流
                  |----> 23bit 坐标    --CDC--> gimbal_ctrl --> pwm_servo
                  |----> frame_tick    --pulse_sync------^
\end{Verilog}

三个角色，三个时钟域：检测/OSD 用像素时钟，云台/PWM 用 50MHz。
跨域只走两座桥（pulse\_sync、cdc\_4phase，都是 Step1 的存货）。

\subsection{做完是什么效果}

屏幕左上角出现状态行：

\begin{Verbatim}
TRK0 LOCK X0560Y0320        <- 跟踪组0 / 已锁定 / 目标坐标
D02 B+00 C064 M0000         <- 延迟帧数 / 亮度 / 对比度 / 运动像素数
2026 FPGA CONTEST - ANLOGIC ...  <- 底部跑马灯
画面里：目标套红色框，中心绿色十字准星，运动目标还有橙色框
\end{Verbatim}

\begin{okbox}{学完 Step4 你能做到}
\begin{itemize}
  \item 解释点阵字库的原理（8x16 ROM，一个字节一行）
  \item 用组合逻辑写出文字区/准星/框的叠加判决
  \item 把二进制坐标变成屏幕上的十进制字符
  \item 说清目标包跨域为什么要"影子寄存器 + 帧首换装"
  \item 用"逐像素金标准比对"证明整帧画面是对的
\end{itemize}
\end{okbox}

% ============================================================
\section{阶段 0：点阵文字的原理（1～2 天）}

\subsection{字就是一个位图}

8x16 点阵字库：每个字符 8 像素宽、16 行高，每行一个字节。
比如字母 A 的第 3 行可能是 \code{8'h3C}（二进制 00111100）——
中间 4 个像素亮。整个字库存成 ROM，地址是
\code{\{字符 ASCII 码, 行号\}}，读出 8 bit 就是这一行的 8 个像素。

\begin{Verilog}
addr = {char, row}          // 12bit: 8bit 字符码 + 4bit 行号
bits = font_rom[addr]       // 8bit, MSB = 最左像素
on   = bits[7 - (px % 8)]   // 当前像素在字节里的位置
\end{Verilog}

\begin{keybox}{为什么字库是"生成"的}
手写 40 个字符 x 16 行的 case 表既痛苦又容易错。工程里字库由
Python 从 OpenCV 的 Hershey 字体渲染生成（\code{gen\_vectors.py}），
可复现、可重新生成——\textbf{生成文件禁止手改}，和 Step2 的倒数表同一条纪律。
\end{keybox}

\subsection{练习}

\begin{enumerate}
  \item 在方格纸上手画 8x16 的 "L" 点阵，写出 16 个字节的十六进制
  \item 思考：MSB=最左像素，如果搞反了会看到什么？
  \item 算一下：1280x720 画面里，16 像素高的文字最多能放几行不重叠？
\end{enumerate}

% ============================================================
\section{阶段 1：坐标系统与叠加判决（2 天）}

\subsection{自己数像素：x/y 计数器}

OSD 模块收到的是 1 像素/拍的流，自己维护 \code{px/py} 计数器就知道
当前画到哪了：

\begin{Verilog}
if (I_valid) begin
    if (I_sof)      px <= 1;      py <= 0;      // 帧首
    else if (I_eol) px <= 0;      py <= py + 1; // 行尾换行
    else            px <= px + 1;
end
\end{Verilog}

valid 门控意味着：流里有消隐气泡（valid 拉低）时计数器自然停走，
\textbf{不需要任何行场同步变换电路}——这是 vpix 流自带 sof/eol 的好处。

\subsection{叠加判决：一层"该不该变色"的组合逻辑}

每个像素问四个问题，按优先级决定颜色：

\begin{Verilog}
文字区:  py 在 [16,32) 且 px 在 [16,176)   -> 查字库, 点亮点用白色
准星:    到 (cx,cy) 的横竖臂, 臂宽 3..12   -> 绿色
目标框:  px 是 minx/minx+1/maxx-1/maxx 且 py 在框内(或对称) -> 红色
运动框:  同上, 用运动 bbox, mv=1 时生效     -> 橙色
都不是:  原样输出底图
\end{Verilog}

\begin{keybox}{准星为什么要"中空"}
十字臂取 \code{|dx| 在 [3,12]} 而不是 [0,12]——中心 3 像素留空，
否则准星正中把目标中心盖住，人眼反而看不清锁定点。这种细节
在金标准渲染器里和 RTL 逐位对齐，两边一致才算对。
\end{keybox}

\subsection{优先级即接线}

实现上就是一个嵌套三元运算符（优先级从上到下）：

\begin{Verilog}
O_data <= on_txt  ? WHITE  :
          on_cross ? GREEN  :
          on_box   ? RED    :
          on_mbox  ? ORANGE : I_data;
\end{Verilog}

\subsection{练习}

\begin{enumerate}
  \item 写出"目标框只画 2 像素厚"的 x 边缘条件（提示：minx 和 minx+1）
  \item 手推：目标在 (100,100)，哪些像素会被画成准星臂？
  \item 思考：为什么判决全部用组合逻辑、只在输出寄存一拍？
\end{enumerate}

% ============================================================
\section{阶段 2：把数字写到屏幕上（2～3 天）}

\subsection{二进制转 ASCII}

坐标寄存器里是 11 bit 二进制数（比如 560），屏幕要显示字符 '5''6''0'
（ASCII 53、54、48）。转换就是反复除 10 取余：

\begin{Verilog}
function [31:0] to_digits4(input [13:0] val);
    d3 = 48 + val / 1000;   r = val % 1000;   // 千位
    d2 = 48 + r / 100;      r = r % 100;      // 百位
    d1 = 48 + r / 10;                          // 十位
    d0 = 48 + r % 10;                          // 个位
endfunction
\end{Verilog}

48 是字符 '0' 的 ASCII 码。\textbf{为什么敢用除法}：这段逻辑每帧只在
"包到达"时跑一次，不在像素流水线上——时序充裕，面积换省事完全正确。

\subsection{拼一条状态行}

20 字符缓冲在包到达时一次性写满：

\begin{Verbatim}
位置:  0  1  2  3  4  5  6  7  8  9  10 11 12 13 14 15 16 17 18 19
内容:  T  R  K  0  (空) L  O  C  K  (空) X  0  5  6  0  Y  0  3  2  0
       ^^^^^ 算法名    ^^^^^ 锁定态    ^^^^ cx 四位十进制 ^^^^ cy
\end{Verbatim}

锁定但没目标时坐标四位替换成 '----'（0x2D），一眼可辨。

\subsection{第二行遥测与跑马灯}

\begin{itemize}
  \item \code{D02}：延迟可视化——用"当前显示帧号 - 包里带的采集帧号"算出
    系统延迟帧数（进阶要求（3））
  \item \code{B+00 C064}：亮度/对比度当前值（扩展要求（4)）
  \item \code{M0000}：运动检测命中像素数（进阶要求（4)）
  \item 跑马灯：64 字符消息，每帧右移 2 像素，取模循环——纯计数器 + ROM
\end{itemize}

\subsection{练习}

\begin{enumerate}
  \item 手推 to\_digits4(560) 每一步的商和余
  \item 如果画面 60fps，跑马灯每帧 2px，扫完整个消息一圈要几秒？
  \item 思考：为什么坐标显示固定 4 位（补零）而不是变长？
\end{enumerate}

% ============================================================
\section{阶段 3：跨时钟域拿目标包（2 天）}

\subsection{问题：检测结果"随时"到达}

检测链在\textbf{帧末消隐期}吐出目标包（135bit），OSD 在\textbf{像素时钟}域
逐拍画图。两边时钟相同（本工程共用像素时钟），但包到达时刻可能在
\textbf{一帧的中间}——直接更新绘制寄存器，会出什么事故？

\begin{warnbox}{撕裂：一帧里上半用旧坐标、下半用新坐标}
目标在第 300 行时包到达，绘制寄存器立即换新值——第 0～299 行画的是
旧位置的框，第 300～719 行画新位置的框。画面上出现\textbf{两个半截框}。
这就是"撕裂"，Step4 v3 专治它。
\end{warnbox}

\subsection{解法：影子寄存器 + 帧首换装}

双缓冲思想：包随时写\textbf{影子}寄存器组；绘制寄存器只在 \code{I\_sof}
拍（新帧第一个像素）从影子\textbf{整组换装}：

\begin{Verilog}
if (pack_valid) begin          // 随时到, 只写影子
    s_cx <= pack[...]; s_cy <= pack[...]; ...
end
if (I_valid && I_sof) begin    // 帧首一次换装
    cx <= s_cx; cy <= s_cy; ...
end
\end{Verilog}

整帧画图期间绘制寄存器\textbf{恒定}——撕裂在结构上不可能发生。
代价是结果最多滞后一帧显示，而 OSD 本来就是帧率更新，无感。

\subsection{135bit 目标包位图（接口唯一权威定义在 IF-VISION.md）}

\begin{Verilog}
[134]    valid          [65:58]  fid (采集帧号)
[133]    algo (0跟踪1边缘) [57]    mot_valid
[132:130] preset 阈值组   [56:43]  mot_cnt
[129:119] cx             [42:32]  mot_minx
[118:108] cy             [31:21]  mot_maxx
[107:97]  minx           [20:11]  mot_miny
[96:86]   maxx           [10:1]   mot_maxy
[85:76]   miny           [0]      保留(必须填 0)
[75:66]   maxy
\end{Verilog}

\begin{warnbox}{bit0 保留位不是装饰}
包拼接时少填一位，整包上移一位，OSD 读 valid 恒 0——画面永远显示
LOST。本项目真出过这个事故（详见 Part II 复盘）。规定：
\textbf{发送端必须把保留位补齐，位偏移只在接口文档定义一处}。
\end{warnbox}

\subsection{练习}

\begin{enumerate}
  \item 手算：cx=560, cy=320，在包里各落在哪些 bit？
  \item 若把 cx 从 [129:119] 错解成 [130:120]，读出的 cx 会偏大还是偏小？
  \item 思考：为什么换装选 sof 而不是 vsync 信号？
\end{enumerate}

% ============================================================
\section{阶段 4：整机集成（2～3 天）}

\subsection{vision\_top：把五个模块拼起来}

\begin{itemize}
  \item \textbf{直通优先}：96bit 入口原样出（O\_tdata=I\_tdata），
    显示链路行为与官方例程完全一致；
  \item \textbf{检测分支}：detect\_top 拿到坐标后拼 135bit 包；
  \item \textbf{云台分支}：23bit 坐标包 + 帧拍分别跨到 50MHz 域；
  \item \textbf{OSD 分支}：串化后的 1ppc 流 + 目标包；
  \item \textbf{声光}：led\_buzzer（锁定绿灯/丢失红灯闪/锁定沿蜂鸣）。
\end{itemize}

\subsection{"包与帧拍异序"的处理}

跨到 50MHz 域后，\textbf{坐标包和帧拍到达顺序不保证}（两个独立 CDC）。
如果帧拍先到而包还没到，云台会用到旧坐标——可接受（1 帧滞后）；
但\textbf{包先到、帧拍后到}时，包不能丢。工程用"缓存 + 帧拍消费"：

\begin{Verilog}
if (ctrl_valid_dst) begin      // 包到了: 缓存, 记下 valid
    ctrl_hold     <= ctrl_pack_dst;
    ctrl_vld_hold <= ctrl_pack_dst[22];
end
if (tick_pwm)                  // 帧拍到: 消费一次, 清 valid
    ctrl_vld_hold <= 1'b0;
\end{Verilog}

gimbal\_ctrl 每帧只看一次 \code{ctrl\_vld\_hold}——无论谁先到，
消费时拿到的都是完整的一套。

\subsection{练习}

\begin{enumerate}
  \item 画出 vision\_top 里三个时钟域各住着哪些模块
  \item 思考：直通输出如果插了一拍延迟，显示会怎样？（提示：和官方链路比对）
  \item ctrl\_hold 为什么不用双缓冲？
\end{enumerate}

% ============================================================
\section{阶段 5：怎么证明画面是对的（1～2 天）}

\subsection{逐像素金标准比对}

"看起来有框"不算数。Python 端有一个 \code{render\_osd()}——
和 RTL 绘制规则\textbf{逐位对齐}的纯软件渲染器（同样知道字库、框内缩、
十字中空）。验证方法：

\begin{Verilog}
RTL 跑 5 帧 -> 整帧转储 (921600 像素/帧, 24bit)
Python 渲染同 5 帧 -> 逐像素比对 -> 全部一致才算 PASS
\end{Verilog}

\begin{keybox}{为什么这个方法是质量基线}
肉眼只能看"有没有框"，看不出"框左移了 1 像素"或"第 37 行少画了一点"。
逐像素比对连\textbf{一个像素的偏差}都逃不掉——Step4 联调抓出的全部 3 个
真缺陷（包宽、旧位、字段溢出）都是它定位的，没有一处靠肉眼。
\end{keybox}

\subsection{练习}

\begin{enumerate}
  \item 算一算：5 帧 x 921600 像素 x 24bit 的转储文件多大？
  \item 思考：如果字库生成器改了一个字符的点阵，哪两边要同步重新生成？
\end{enumerate}

% ============================================================
\section{阶段 6：跑通测试（1～2 天）}

\subsection{一键回归与怎么读结果}

\begin{Verbatim}
cd D:\fpga-ptz-tracker\tools
python run_all.py tb_osd_track tb_vision_top tb_closed_loop
\end{Verbatim}

\begin{Verbatim}
--- tb_osd_track ...
    # FRAME0 OK hash=一致   <- 整帧 921600 像素逐位一致
    ==> PASS

--- tb_vision_top ...
    # PASSTHRU OK          <- 直通逐拍一致
    # FRAME4 OK            <- OSD 帧与金标准一致
    # DUTY x=1428 y=1500   <- 云台朝目标方向动, Y 轴死区正确
    ==> PASS

--- tb_closed_loop ...
    TB_RESULT: PASS settle_frame=12 final_err=-8
    ==> PASS               <- 闭环 12 帧锁定
\end{Verbatim}

\begin{okbox}{Step4 学习清单（全部打勾就过关）}
\begin{itemize}
  \item 能画 8x16 字库的取址与点亮过程
  \item 能写出文字/准星/框的判决条件并解释优先级
  \item 能解释影子寄存器为什么能消撕裂
  \item 能背出 135bit 包的字段顺序，说清保留位的教训
  \item 能解释"逐像素金标准比对"为什么比肉眼强
  \item 跑通 3 个测试并读懂输出
\end{itemize}
\end{okbox}

% ============================================================
% Step4 教程 · Part II：工程实现详解
% 由 step4_tutorial.tex \input 引入（共用同一份导言区与主题色）
% ============================================================

\part{工程实现详解}

\begin{tipbox}{这一部分写什么}
仓库里真实做出来的 Step4 工程（RTL 共 1231 行：字库 684 + osd\_track 313 +
vision\_top 234）逐模块讲透，外加三次联调迭代的完整复盘。
对应目录：\code{03\_vision/step4\_osd\_integration/}。
\end{tipbox}

% ============================================================
\section{工程总览}

\subsection{文件清单}

\begin{center}
\begin{tabular}{llr l}
\toprule
\rowcolor{rowtint}\textbf{文件} & \textbf{类型} & \textbf{行数} & \textbf{职责} \\
\midrule
\code{font\_rom\_8x16.v} & 生成 & 684 & 40 字符 ASCII 子集点阵（脚本生成） \\
\code{osd\_track.v} & RTL & 313 & OSD 叠加层 v3（三行文本+双框+双缓冲） \\
\code{vision\_top.v} & RTL & 234 & 整机顶层 v2（检测+OSD+云台+声光） \\
\code{tb/tb\_osd\_track.v} & 测试 & 126 & OSD 整帧逐像素比对 \\
\code{tb/tb\_vision\_top.v} & 测试 & 99 & 整机 5 断言 \\
\code{tb/tb\_closed\_loop.v} & 测试 & 95 & 闭环在环收敛 \\
\code{sim\_out/*.png} & 证据 & --- & OSD 渲染对比图（金标准 vs RTL） \\
\bottomrule
\end{tabular}
\end{center}

\subsection{三时钟域与两条跨域河}

\begin{Verilog}
pix 域 (74.25MHz):  检测链 / OSD / 显示路径
pwm 域 (50MHz):     gimbal_ctrl / pwm_servo
跨域信号:
   frame_tick    pulse_sync (单 bit)
   23bit 坐标包  cdc_4phase
   135bit 目标包 cdc_4phase (同域备用, 接口兼容异域)
\end{Verilog}

\begin{keybox}{一条贯穿全项目的设计规则}
任何跨域信号只能走 pulse\_sync 或 cdc\_4phase 两座桥，禁止手搓两级同步器。
本步的 135bit 包就是 \code{cdc\_4phase \#.WIDTH(135)} 的一次参数化复用——
Step1 写好的桥，Step4 直接搬。
\end{keybox}

% ============================================================
\section{模块详解：font\_rom\_8x16（生成文件）}

字库 ROM 覆盖 OSD 用到的 40 个 ASCII 字符（大写字母、数字、空格、
\code{- + \% . :} 等）。生成链路：

\begin{Verilog}
cv2.putText(Hershey 字体) -> 8x16 阈值化点阵 -> case 表 RTL
地址 {char[7:0], row[3:0]} -> 该行 8bit (MSB = 最左像素)
\end{Verilog}

为什么选 8x16：16 像素高在 720p 上恰好是 45 行文字的量级，
评委在演示视频里看得清；8 像素宽让字符地址 = 像素 x 右移 3 位，
判决逻辑零乘法。

% ============================================================
\section{模块详解：osd\_track v3}

\subsection{v1 到 v3 的演进}

\begin{center}
\begin{tabular}{lp{8.6cm}}
\toprule
\rowcolor{rowtint}\textbf{版本} & \textbf{内容} \\
\midrule
v1 & 单行文字 + 准星 + 单框，68bit 包（IF-VISION v1.0.0） \\
v2 & 配合 Step5 拓展：包扩到 135bit，三行文本 + 跑马灯 + 运动框 \\
v3 & 消撕裂：影子寄存器组 + sof 拍整组换装（vsync 语义） \\
\bottomrule
\end{tabular}
\end{center}

\subsection{影子寄存器组：撕裂的结构性根治}

v3 的核心是\textbf{两组一模一样的寄存器}：

\begin{itemize}
  \item \textbf{影子组}（\code{s\_*}）：\code{pack\_valid} 拍写入，
    时刻可能在帧中间；
  \item \textbf{绘制组}（无前缀）：只在 \code{I\_valid \&\& I\_sof} 拍
    从影子整组拷贝，帧内再也不会变。
\end{itemize}

绘制判决只读绘制组——无论包何时到，一帧之内的坐标/文字/框
\textbf{是同一份快照}。已知残留：sof 拍本身（像素 (0,0)）仍用旧值绘制
一拍，但 OSD 起始在 x=16，无视觉影响（RTL 注释里明确记了这笔账）。

\subsection{文本构建：包到达拍的一次性算术}

\begin{Verilog}
txt1_s[0..2] = pack[133] ? "EDG" : "TRK"      // 算法名
txt1_s[3]    = 48 + preset                     // 阈值组号
txt1_s[5..8] = pack[134] ? "LOCK" : "LOST"    // 锁定态
txt1_s[11..14] = to_digits4(cx) 或 "----"      // 坐标四位
txt2: "D" lat " B" ±bb " C" cc " M" mmmm       // 遥测行
\end{Verilog}

延迟显示 \code{lat = disp\_cnt - fid}：包里带的是\textbf{采集时刻}的帧号，
显示端用当前帧号一减，就是这套链路的真实延迟帧数（进阶（3）的可视化）。

\subsection{跑马灯的算术}

\begin{Verilog}
mar_off = disp_cnt * 2                  // 每帧 2px
mar_pos = px - TXT_X0 + mar_off         // 相对位置 + 滚动偏移
mar_ci  = mar_pos[8:3] % 63             // 字符索引(63 字符取模)
on_mar  = 字形位[7 - mar_pos[2:0]]      // 字节内位选
\end{Verilog}

63 字符 x 8 像素 = 504 像素周期，2px/帧 即 252 帧（约 4.2 秒）滚一圈。

\subsection{准星与双框：带符号差 + 边缘四线}

\begin{Verilog}
dx_c = $signed(px) - $signed(cx);   // 带符号差, 一切判决的基础
ax = (|dx_c| <= 12) && (|dx_c| >= 3);   // 臂: 3..12 (中心 3px 留空)
on_cross = tv && ((ay && dx_c==0) || (ax && dy_c==0));
// 框: 四条边各 2px 厚
x_edge = (px==minx) || (px==minx+1) || (px==maxx-1) || (px==maxx);
on_box = tv && ((x_edge && in_by) || (y_edge && in_bx));
\end{Verilog}

\begin{warnbox}{带符号差是框类判决的唯一正确姿势}
如果用无符号减法算 \code{px - cx}，目标左半边的像素差全部回绕成
巨大正数，左半边框直接消失——和 Step3 闭环跑飞是同一类根因。
本项目已两次被"无符号减法"咬，写进了方法论。
\end{warnbox}

\subsection{输出 mux 与颜色}

优先级一个表达式写完（文字 $>$ 准星 $>$ 目标框 $>$ 运动框 $>$ 底图）：

\begin{Verilog}
O_data <= on_txt1|on_txt2 ? 24'hFFFFFF :   // 文字 白
          on_mar   ? 24'h00FFFF :          // 跑马灯 青
          on_cross ? 24'h00FF00 :          // 准星 绿
          on_box   ? 24'hFF0000 :          // 目标框 红
          on_mbox  ? 24'hFF8000 : I_data;  // 运动框 橙
\end{Verilog}

% ============================================================
\section{模块详解：vision\_top v2}

\subsection{算法切换：帧边界平稳换源}

进阶（1）要求多算法切换。工程做法：\textbf{边缘检测常开}（并行分支），
切换只动一个 mux 的选择端，且\textbf{在帧首拍锁存}：

\begin{Verilog}
if (I_tvalid && I_tuser)          // 只在帧首采样
    algo_q <= I_algo_edge;
// 之后一整帧 mux 固定选检测流或边缘流
\end{Verilog}

帧中间切会怎样？mux 前两条流水线延迟不同，切换拍前后像素来自
\textbf{两条链的不同进度}，画面错位——帧边界锁存让切换永远发生在
"两帧之间"，视觉无缝。

\subsection{135bit 包拼装与保留位}

\begin{Verilog}
wire [134:0] pack_vision = {t_valid, I_algo_edge, I_preset_idx,
    t_cx, t_cy, t_minx, t_maxx, t_miny, t_maxy, t_fid,
    mo_valid, mo_cnt[13:0], mo_minx, mo_maxx, mo_miny, mo_maxy,
    1'b0};    // <- bit0 保留位, 必须补 0
\end{Verilog}

RTL 注释里直接写着事故记录："缺位时整包上移一位，OSD 解析 valid 恒 0
（永远 LOST 事故根因，2026-09-20 修）"。\textbf{把教训写在案发现场}，
下一个改包的人第一眼就能看到。

\subsection{控制路径：缓存 + 帧拍消费}

23bit 控制包（valid+cx+cy）跨域后进 \code{ctrl\_hold} 缓存，
\code{ctrl\_vld\_hold} 标记有效；帧拍 \code{tick\_pwm} 到达时消费一次。
这样"包先到/帧拍先到"两种顺序都能正确工作，
代价是控制滞后 1 帧——设计行为，闭环增益公式里已经算过它。

\subsection{声光反馈（给嵌入式方向的接口）}

\code{led\_buzzer}：锁定亮绿灯、丢失红灯 1Hz 闪、\textbf{锁定沿}（从丢到锁）
蜂鸣 200ms。演示时评委不用盯着屏幕，听声音就知道系统锁上了。

% ============================================================
\section{测试平台设计}

\subsection{tb\_osd\_track：整帧逐像素}

激励：\code{osd\_coord.txt} 给出每帧的目标坐标/包围框（金标准写的），
TPG 底图 + OSD 叠加后整帧转储；期望：\code{osd\_exp*.hex} 是
Python \code{render\_osd()} 渲染的同帧画面。比对粒度：\textbf{像素级
24bit 相等}，4 帧共 3686400 个样本。

\subsection{tb\_vision\_top：五条整机断言}

\begin{center}
\begin{tabular}{ll}
\toprule
\rowcolor{rowtint}\textbf{断言} & \textbf{意义} \\
\midrule
直通输出与输入逐拍一致 & 主线零扰动（显示不受影响） \\
第 4 帧 OSD 与金标准一致 & 整链绘制的像素级正确 \\
duty\_x 在 [700,1500) & 云台 X 朝目标方向动 \\
duty\_y == 1500 & Y 轴死区正确（误差 20 = DZ\_Y 边界内） \\
lost\_cnt == 0 & 无误丢帧 \\
\bottomrule
\end{tabular}
\end{center}

\begin{warnbox}{迭代 3：采样过早的教训}
第一版在第 2 帧末读 duty，恒 1500 判 FAIL——其实云台正常，只是
"检测(1帧)+除法+跨域+消费"共约 2 帧滞后是设计行为。
修复：采样点移到第 4 帧末。\textbf{给断言选时机，要按数据通路的
真实延迟账来}，不能拍脑袋。
\end{warnbox}

\subsection{tb\_closed\_loop：相机模型驱动 TPG}

TPG 方块位置由 TB 相机模型实时计算：
\code{box\_x = 640 + (40 度 - theta) x 20px/度}——云台转多少，
方块就往回挪多少，\textbf{和真实相机的行为同构}。60 帧判据：
误差进入 ±25px 死区并保持 10 帧。实测 \code{settle\_frame=12}
（T1 控制节拍修正后；修正前为 36），
逐帧轨迹在 \code{sim/closed\_loop\_trace.txt}（答辩可直接画收敛曲线）。

% ============================================================
\section{验证结果与联调复盘}

\subsection{回归矩阵}

\begin{center}
\begin{tabular}{lll}
\toprule
\rowcolor{rowtint}\textbf{测试} & \textbf{比对粒度} & \textbf{结果} \\
\midrule
tb\_osd\_track & 整帧逐像素（4 帧 3686400 样本） & PASS \\
tb\_vision\_top & 5 断言 + 整帧 & PASS \\
tb\_closed\_loop & 在环收敛（settle\_frame） & PASS：12 帧 \\
\bottomrule
\end{tabular}
\end{center}

\subsection{"永远 LOST"事故复盘（本步骤最大缺陷）}

\begin{warnbox}{一次事故，三个裂缝}
现象：OSD 永远显示 LOST、云台 duty 恒 1500 不动、文字区坐标被污染。\\[3pt]
\textbf{根因链}：vision\_top 拼 69bit - 截成 67bit（valid 位丢）；
TB 侧 67bit 接 68bit 端口（MSB 悬空为 Z）；osd\_track 按 67bit 解包。
\textbf{三处三个宽度}，每一个单看都"能用"。\\[3pt]
\textbf{连带发现}：位偏移重排后，文字分支还读旧位 [66]（现在是 preset 位）
而不是 [67]（valid 位）——所以出现"显示 LOST 但坐标数字正常"
的灵异现象（数字用的是 [67]）。一字符修复。\\[3pt]
\textbf{修复动作}：统一宽度、位偏移重排、\textbf{冻结进 IF-VISION.md}
（接口唯一权威定义），此后任何位宽变更必须走接口变更记录。
\end{warnbox}

\begin{keybox}{为什么逐像素比对是抓这类 bug 的唯一手段}
包宽错误不会让仿真崩溃——valid 恒 0 的表现只是"画面对但没框"，
肉眼完全可能当成"场景里没有目标"。只有逐像素比对才把"应有框 vs 实际无框"
的每个像素差摆上台面。Step4 的三个真缺陷（包宽、旧位、字段溢出）
全部由它定位。
\end{keybox}

\subsection{Step4 真缺陷清单}

\begin{center}
\begin{tabular}{llp{5.6cm}l}
\toprule
\rowcolor{rowtint}\textbf{编号} & \textbf{现象} & \textbf{根因} & \textbf{类型} \\
\midrule
O1 & 永远 LOST & 包宽三处不一致（69/67/68） & 接口 \\
O2 & LOST 但坐标正常 & word 分支读旧位 [66] & RTL \\
O3 & 框画到 (72,18) & 坐标文件 \%02X 写 300 溢出 3 字符 & TB 脚本 \\
O4 & duty 恒 1500 误报 & TB 采样点早于通路真实延迟 & TB \\
\bottomrule
\end{tabular}
\end{center}

% ============================================================
\section{上板挂载说明（给上板组）}

\begin{center}
\begin{tabular}{ll}
\toprule
\rowcolor{rowtint}\textbf{项目} & \textbf{内容} \\
\midrule
挂载方式 A & 显示侧：官方 hdmi\_mixer 之前串入 osd\_track（vpix 域） \\
挂载方式 B & 检测侧：串在检测分支出口（本工程仿真即此形态） \\
vpix/DE 转换 & 显示侧若只有 DE 信号，按 IF-OSD.md 的规则生成 sof/eol \\
按键 & preset\_idx（阈值组）/ track\_en / algo\_edge / 亮度对比度 \\
观测 & ILA 挂 135bit 包 + O\_status；duty\_dbg 看云台 \\
\bottomrule
\end{tabular}
\end{center}

\begin{tipbox}{答辩演示建议}
画面里的三行 OSD 本身就是"评分表逐条自证"：L1 证明扩展（1）叠加与
进阶（2）坐标标注，D 字段证明进阶（3）延迟可视化，M 字段证明
进阶（4）运动检测，B/C 字段证明扩展（4）实时调节——
不需要额外口头解释。
\end{tipbox}

% ============================================================
\section{小结：Step4 的可复用产出}

\begin{itemize}
  \item \textbf{整机顶层 vision\_top}：检测/OSD/云台/声光四合一，
    上板即用；
  \item \textbf{一套 OSD 叠加层}：三行文本、双框、跑马灯、双缓冲，
    全部逐像素验证；
  \item \textbf{逐像素金标准比对方法}：此后任何绘制改动都有回归兜底；
  \item \textbf{接口治理范例}：135bit 包唯一权威定义 + 变更记录制度；
  \item \textbf{收敛证据链}：closed\_loop\_trace.txt 逐帧轨迹，
    答辩直接可画。
\end{itemize}

\begin{okbox}{下一步（Step5 预告）}
基础与进阶都通了，Step5 把评分表剩下的扩展项逐个补齐：
多算法切换、运动检测、亮度对比度、声光反馈——
每个都带位精确验证，并最终把回归扩到 16 个测试全绿。
\end{okbox}


% ============================================================
\part*{结语}
\clearpage
\section*{总结与下一步}
\addcontentsline{toc}{section}{总结与下一步}

\subsection*{Step4 学到了什么}
\addcontentsline{toc}{subsection}{Step4 学到了什么}
\begin{itemize}
  \item 点阵字库原理与生成式字库（cv2 Hershey 渲染）
  \item vpix 流上的 x/y 计数与叠加判决（文字/准星/双框/优先级）
  \item 二进制转 ASCII（常量除法）与 20 字符文本缓冲
  \item 影子寄存器 + 帧首换装：结构上消灭撕裂
  \item 135bit 目标包位图与接口单一权威定义
  \item 整机集成：三时钟域、直通零扰动、包与帧拍异序消费
  \item 逐像素金标准比对方法学
\end{itemize}

\subsection*{下一步学什么}
\addcontentsline{toc}{subsection}{下一步学什么}
\begin{itemize}
  \item \textbf{Step5 拓展功能集}：把评分表上的进阶/扩展项逐个补齐
  \begin{itemize}
    \item 多算法切换（边缘检测并行分支 + 帧边界无缝切换）
    \item 帧差运动检测与橙色包围框
    \item 亮度/对比度实时调节、声光反馈、三行 OSD 与跑马灯
  \end{itemize}
  \item 对应文档：\code{03\_vision/step5\_expansion/}
  \item 对应金标准：\code{tools/golden\_model.py::frame\_diff\_motion}
\end{itemize}

\end{document}
